% ====================================================================
% Section 3 (v2): The two-term physical core.
% Replaces the 8-layer "mass-aware" model of v1. Faithful to
% experiments/emmet_core.py.
% ====================================================================
\section{The Two-Term Physical Core}
\label{sec:core}

\noindent\emph{(Presented as built: Section~\ref{sec:redux} re-audits this core on the modern bench and inverts which term earns its place.)}

We model a packet as a classical particle descending a potential field
defined on the network graph. The particle is subject to two forces, and
only two. We state them first, then justify by ablation
(Section~\ref{sec:ablation}) that no further mechanism is needed.

\subsection{The potential}

Each directed traversal of an edge $(u,v)$ carries a cost
\begin{equation}
\phi(u,v) \;=\; \alpha\,\ell(u,v) \;+\; \beta\,\frac{x(u,v)}{c(u,v)} \;+\; \gamma\,s(u,v),
\label{eq:potential}
\end{equation}
where $\ell$ is edge latency, $x/c$ is the instantaneous load-to-capacity
ratio (congestion), and $s$ is a persistent \emph{scar} term defined below.
The three weights $(\alpha,\beta,\gamma)$ are the only tunable quantities of
the model. All three are local: computing $\phi$ needs no global state beyond
the edge itself.

\subsection{Force 1: gradient descent}

The particle follows the downhill direction of $\phi$ from source to
destination, subject to a hop budget: it may use a path no longer than
$\lceil \alpha_{\text{budget}}\, h_{\text{sp}} \rceil$ hops, where
$h_{\text{sp}}$ is the shortest-path hop count. The minimum-potential path
under this budget is found exactly by a dynamic program over
$(\text{hops used},\, \text{node})$, in $O(k\,|E|)$ time with
$k=\lceil\alpha_{\text{budget}} h_{\text{sp}}\rceil$.
The budget is what lets the particle take a slightly longer, less congested
detour instead of forcing every packet onto the single shortest path---the
mechanism by which load spreads.

% Algorithm 1: the two-term core routing step. Uses algorithm2e.
\begin{algorithm}[t]
\caption{EMMET-core routing step (two-term physical router)}
\label{alg:core}
\SetAlgoLined
\KwIn{graph $G=(V,E)$ with per-edge latency $\ell$, load $x$, capacity $c$;
scar field $s$; source $a$, destination $z$; weights $\alpha,\beta,\gamma$;
budget factor $\alpha_b$}
\KwOut{path $a \rightsquigarrow z$}
$h_{\text{sp}} \leftarrow$ shortest-path hop count $a \to z$\;
$k \leftarrow \lceil \alpha_b \, h_{\text{sp}} \rceil$\tcp*{hop budget}
\For{each edge $(u,v)$}{
  $\phi(u,v) \leftarrow \alpha\,\ell(u,v) + \beta\,x(u,v)/c(u,v) + \gamma\,s(u,v)$\;
}
$P \leftarrow$ min-$\sum\phi$ path of length $\le k$ via DP over (hops, node)\tcp*{$O(k|E|)$}
\Return $P$\;
\BlankLine
\tcp{On a drop at edge $(u,v)$ while forwarding:}
$s(u,v) \leftarrow s(u,v) + 1$\tcp*{Newton-III reaction}
\tcp{Once per step, everywhere:}
$s \leftarrow \rho\, s$, \quad $\rho = 2^{-1/T_{1/2}}$\tcp*{scars fade}
\end{algorithm}


\subsection{Force 2: Newton's third law (the scar)}

The gradient term reacts to congestion that is \emph{present}. It cannot
react to congestion that has just \emph{caused a loss}: once a packet is
dropped, the offending edge's instantaneous load may already have fallen.
We close this gap with a reaction term inspired by Newton's third law---to
every action a packet undergoes, the network responds with an equal and
opposite reaction on the packets that follow.

Concretely, when a packet is dropped on edge $(u,v)$, that edge's scar is
incremented, $s(u,v) \leftarrow s(u,v) + 1$. The scar enters the
potential~\eqref{eq:potential} with weight $\gamma$, so subsequent particles
feel a repulsive reaction steering them away from edges that recently
failed---information the instantaneous congestion term does not carry.

Scars are not permanent: at each step they decay,
$s(u,v) \leftarrow \rho\, s(u,v)$ with $\rho = 2^{-1/T_{1/2}}$ and half-life
$T_{1/2}$. A corridor that stops failing reopens on its own, and at rest
(no losses) the scar field vanishes and the router reduces to pure gradient
descent. The mechanism is self-regulating: it acts only when and where the
network is actually breaking. This is the only memory in the model, and---as
the ablation will show---the only added mechanism that earns its keep.

\subsection{What is deliberately absent}
\label{sec:absent}

This two-term core is the survivor of a larger model. Earlier versions of
EMMET endowed the packet-particle with extra physical attributes: an
effective mass that grew with traversed congestion, a network thermostat
modulating the congestion weight, a temporal ripple diffusing loss memory
over several steps, and an Archimedes-style buoyancy term. Each was
physically motivated and individually plausible.

A systematic ablation (Section~\ref{sec:ablation}) removed them one at a time
and measured the cost. The result was unambiguous: the gradient and the
Newton-III scar account for essentially all of the performance, while the
remaining mechanisms ranged from negligible to mildly harmful---removing the
ripple and the thermostat slightly \emph{improved} results. We therefore
report the two-term core as the model and treat the richer apparatus as a
cautionary tale about physically-seductive terms that do not survive
measurement. The reference implementation is 60 lines of Python.

